\section{Introduction}

A system that exists at time \(t\) has survived until time \(t\). Much less follows from that statement than is often assumed.

Its present condition may have been inherited from earlier activity rather than generated by current processes. A currently functioning system may be consuming reserves faster than it replaces them. A participant that appears useful in isolation may reduce collective viability once network responses are considered. A regulator that is individually favoured may provide less control than the wider system requires. A signal used to reinforce a process may gradually cease to measure the function for which the process was originally retained.

These are not unusual exceptions. They arise whenever the processes responsible for maintaining an organization are distributed across interacting components with different incentives, time scales, dependencies, signals, and failure modes.

The relevant question is therefore not merely: \emph{What state is the system currently in?} Nor even: \emph{Which components contribute to that state?} The stronger question is:

\begin{quote}
\textbf{Which current processes maintain the capacity of the system to remain within its declared viability conditions under future disturbance, and what maintains those processes themselves?}
\end{quote}

This question connects several mature traditions.

Viability theory studies whether dynamical systems can remain within specified constraint sets and characterizes the states and controls from which continued viability is possible \citep{aubin1991,aubin2011}. Organizational approaches to biology describe organisms as networks of mutually dependent constraints that participate in one another's production and maintenance \citep{mossio2010,montevil2015,moreno2015}. Di Paolo's account of adaptivity links regulation to conditions of viability and organism-relative normativity \citep{dipaolo2005}. Bich and colleagues characterize biological regulation as a form of second-order control that modulates constitutive processes in response to perturbation \citep{bich2016}; subsequent work emphasizes heterarchically organized control rather than necessarily centralized control \citep{bich2022}.

The present paper does not replace these theories. Its narrower purpose is to connect them to a recurring problem exposed by a sequence of formal and computational models of distributed regulation: \emph{the variables that describe current organization are not sufficient to describe the processes that maintain the future possibility of that organization}.

The synthesis developed here is built around boundary conditions rather than a claim of successive scientific falsifications. Each modelling exercise exposed a point at which the previous representation ceased to be sufficient. Several resulting mathematical statements are elementary, and several recover known principles from reliability engineering, control theory, cooperation research, viability theory, and adaptive-network theory. Their value here is compositional: taken together, they suggest a diagnostic framework for distributed maintenance.

The central proposal is
\begin{center}
\fbox{\parbox{0.88\linewidth}{\centering
Distributed maintenance concerns the preservation of viability-relevant
relations and capacities, not merely the persistence of present state.}}
\end{center}

This leads to six separations:
\[
\begin{aligned}
\text{local selection} &\not\Rightarrow \text{sufficient control},\\
\text{persistence} &\not\Rightarrow \text{present self-maintenance},\\
\text{presence} &\not\Rightarrow \text{positive contribution},\\
\text{state} &\not\Rightarrow \text{corrective capacity},\\
\text{reinforcement signal} &\not\Rightarrow \text{viability-relevant function},\\
\text{one positive contribution} &\not\Rightarrow \text{global benefit}.
\end{aligned}
\]

The paper asks what survives once these implications are removed.

\section{Relation to the preceding programme}

The framework developed here should not be read as re-presenting several results already developed in companion work.

\subsection{Persistence does not measure function}

\emph{Persistence Does Not Measure Function} separates equilibrium persistence, productive efficiency, maintained mass, structural support, and externally specified functional capacity \citep{vanhoek2026persistence}. Its counterexamples show that a system can improve according to one of these quantities while deteriorating according to another.

The present paper takes that separation as given. It asks what additional variables are needed once the object of study becomes the \emph{maintenance of future viability} rather than present function alone.

\subsection{Selected regulation does not guarantee sufficient regulation}

\emph{When Does Regulation Pay?} distinguishes the controller strength selected under a stated endogenous objective from the strength required to satisfy an independently declared functional condition \citep{vanhoek2026regulation}. In its minimal model,
\[
k_{\rm opt}\neq k_{\rm func}
\]
in general, and regulation may be retained while remaining too weak to satisfy the functional boundary.

The present paper does not reclaim this as a new result. Instead it generalizes the inferential lesson:
\[
\boxed{
\text{the process selected locally need not supply the quantity required collectively.}
}
\]

\subsection{Present persistence does not establish present maintenance}

A separate test programme distinguishes current stock from the current processes replacing loss and restoring disturbances \citep{vanhoek2026present}. A system can therefore appear healthy while inherited capacity is being depleted.

Again, the present paper takes this result as a starting condition. Its added question is how such maintenance is distributed across a network, how individual contributions depend on architecture and state, and how the maintenance processes themselves persist.

Thus the present manuscript is not a replacement for these companion papers. It is a synthesis of the boundary conditions that become visible when their distinctions are applied to decentralized organization.

\section{Claim-status map}

Because this manuscript combines prior art, results established elsewhere in the author's programme, exact consequences of declared toy models, and new synthesis, those categories are kept explicit throughout. \Cref{tab:claims} summarizes the status of the main claims.

\begin{table}[htbp]
\centering
\small
\caption{Claim-status map. ``Exact model result'' means exact within the stated model; it does not imply empirical validation or literature priority.}
\label{tab:claims}
\begin{tabularx}{\textwidth}{@{}p{0.05\textwidth}X p{0.21\textwidth}@{}}
\toprule
ID & Claim & Status \\
\midrule
S1 & Persistence, productive efficiency, maintained mass, structural support, and functional capacity need not coincide. & \claimstatus{previous paper} \\
S2 & Locally selected regulatory strength need not be sufficient for an independently declared functional boundary. & \claimstatus{previous paper} \\
S3 & Present state does not in general identify present corrective or renewal capacity. & \claimstatus{previous work + exact model} \\
S4 & Participant contribution is relational: its sign can change after network response or across viability criteria. & \claimstatus{exact model / established analogue} \\
S5 & In sampled or capacity-limited correction, participant addition can alter causal access to correction; all-participant monotonicity does not automatically transfer. & \claimstatus{JAM companion paper} \\
S6 & Mixed-sign effects across several viability margins have no weight-independent global scalar sign. & \claimstatus{elementary exact result} \\
S7 & Corrective capacity that decays or is consumed must itself receive sufficient renewal if it is to persist. & \claimstatus{elementary exact result / prior art} \\
S8 & A signal used to allocate renewal can diverge from the viability-relevant function it is meant to track. & \claimstatus{new synthesis / hypothesis} \\
S9 & ``Function--renewal proxy gap'' is proposed as a diagnostic variable for that divergence; no universal metric is claimed. & \claimstatus{new terminology} \\
S10 & A distributed maintenance architecture can be represented as viability specification, informative coupling, response, independent corrective capacity, return, capacity renewal, and architectural revision. & \claimstatus{proposed synthesis} \\
S11 & Viability-conditioned requirements have the form model + declared viability property \(\Rightarrow\) admissible region; they are not categorical moral obligations. & \claimstatus{viability-theory interpretation} \\
S12 & The framework earns empirical value only if its explicit maintenance variables improve prediction or intervention beyond state-only or trend-only models. & \claimstatus{test criterion} \\
\bottomrule
\end{tabularx}
\end{table}

\section{A minimal language of distributed viability}

Let the system state at time \(t\) be \(z_t\in Z\). Let its dynamics be
\[
\boxed{
z_{t+1}=F(z_t,\theta,d_t),
}
\]
where \(\theta\) describes architecture, control parameters or policies; \(d_t\in\D\) is a disturbance; and \(\D\) is the declared disturbance class.

Let
\[
\V\subseteq Z
\]
denote a declared viability region. The one-step admissible region is
\[
\boxed{
\A_{\V}(z)
=
\left\{
\theta:
F(z,\theta,d)\in\V
\quad
\forall d\in\D
\right\}.
}
\]
Longer-horizon viability asks whether some admissible policy keeps the resulting trajectory within \(\V\). This is standard territory for viability and invariance theory.

The new terminology introduced here concerns the architecture implementing those admissible dynamics.

\paragraph{Definition (distributed maintenance architecture).}
A \emph{distributed maintenance architecture} is a collection of local processes whose state-dependent actions jointly alter the probability or margin by which one or more declared viability conditions will continue to hold.

This definition deliberately does not require autonomous agents, explicit rewards, conscious goals, centralized monitoring, biological individuality, or a single scalar system objective. The word \emph{maintenance} therefore refers to a functional relation between present processes and future viability, not to a particular mechanism.

\section{From state to margin}

Membership in \(\V\) is often too coarse for studying maintenance. We therefore introduce a viability margin
\[
M(z,\theta),
\]
where \(M>0\) denotes remaining room before a declared boundary is crossed, \(M=0\) the boundary itself, and \(M<0\) failure of that particular criterion.

The exact definition is application-specific. For a simple stock \(K\) with viability threshold \(V\),
\[
M=K-V.
\]
For a distributed verification system it might be a log safety-certificate margin. For a population it might be distance from a transmission or extinction boundary.

For several simultaneous requirements, define
\[
\boxed{
\Mvec=(M_1,\ldots,M_k).
}
\]
This apparently modest change matters because a distributed process can affect different viability margins in opposite directions. The system therefore need not admit a natural scalar statement of ``better.''
